\documentclass[11pt]{article}
\newcommand{\coursedocname}{Problem Set 0}

\usepackage[utf8]{inputenc}
\usepackage[T1]{fontenc}
\usepackage{lmodern}
\usepackage[margin=1in]{geometry}
\usepackage[dvipsnames]{xcolor}
\usepackage{graphicx, booktabs, enumitem, float, framed, caption, multirow}
\usepackage{fancyhdr}
\usepackage{tcolorbox}
\usepackage{listings}
\usepackage{tikz}
\usepackage{setspace}
\usepackage{needspace}
\IfFileExists{algorithm.sty}{%
  \usepackage{algorithm}
  \usepackage[noend]{algpseudocode}
  \newcommand{\hasalgorithms}{}%
}{}
\usepackage{amsmath,amssymb,bm,xcolor}




\newcommand{\lr}{\alpha}                %

\newcommand{\E}{\mathbb{E}}             %
\newcommand{\Var}{\mathrm{Var}}        %

\usepackage{hyperref}

\definecolor{DartmouthGreen}{HTML}{00693E}
\newcommand{\dgreen}[1]{\textcolor{DartmouthGreen}{#1}}
\newcommand{\dgbf}[1]{\textbf{\dgreen{#1}}}      %
\urlstyle{same}
\hypersetup{
    colorlinks=true,
    linkcolor=DartmouthGreen,
    filecolor=DartmouthGreen,
    citecolor=DartmouthGreen,
    urlcolor=DartmouthGreen,
    breaklinks=true
}

\newcommand{\coursenumber}{COSC 81/281}
\newcommand{\courseterm}{Fall 2026}
\newcommand{\gradescopelink}{\textbf{Gradescope} (\url{https://www.gradescope.com/courses/1372528})}

\newcommand{\psZeroDue}{Tuesday, September 22 (11:59pm ET)}

\setlength{\parindent}{0pt}
\setlength{\parskip}{6pt plus 2pt minus 1pt}
\setlength{\headheight}{13.6pt}
\setlength{\emergencystretch}{3em}
\providecommand{\tightlist}{\setlength{\itemsep}{0pt}\setlength{\parskip}{0pt}}
\setcounter{secnumdepth}{0}

\DeclareRobustCommand{\callout}[2][gray!10]{%
  \begin{tcolorbox}[left=8pt, arc=0pt, outer arc=0pt,
                    colframe=#1, colback=#1, coltext=black]#2\end{tcolorbox}}




\newcommand{\gradstar}{$\bigstar$~\textsf{\small (required for 281; optional for 81)}~}
\newcommand{\bonusq}{\textsf{\small\textbf{[Bonus]}}~}
\newcommand{\bonus}{\bonusq}

\newcommand{\setupheader}[1]{%
  \fancypagestyle{default}{%
      \lhead{#1}
      \chead{}
      \rhead{\textbf{\coursenumber}}
      \lfoot{}
      \cfoot{}
      \rfoot{\thepage}
      \renewcommand{\headrulewidth}{0.4pt}
      \renewcommand{\footrulewidth}{0.4pt}
  }
  \pagestyle{default}
  \fancypagestyle{plain}{\pagestyle{default}}
}
\newcommand{\setupexamheader}[1]{%
  \fancypagestyle{default}{%
      \lhead{#1}
      \chead{Name:}
      \rhead{}
      \lfoot{}
      \cfoot{}
      \rfoot{\thepage}
      \renewcommand{\headrulewidth}{0.4pt}
      \renewcommand{\footrulewidth}{0.4pt}
  }
  \pagestyle{default}
  \fancypagestyle{plain}{\pagestyle{default}}
}
\newcommand{\pstitle}[2]{%
  \title{\textbf{\coursenumber:} #1}
  \author{\textbf{Due} #2\\
  }
  \date{}
}

\newcommand{\coursenotice}{%
  \callout{Problem sets are graded \textbf{pass/fail on completion}, and
  there are no late days. The problem sets are where the learning happens,
  and they are the best available practice for the midterm and final exam.
  Submit to \gradescopelink\ and assign each question to a page in your
  submission. Questions are welcome on \dgbf{Slack}. Starred ($\bigstar$)
  questions are required for COSC 281 students and optional (ungraded) for
  COSC 81 students.}}

\lstset{
  language=python,
  basicstyle=\ttfamily,
  commentstyle=\color{BrickRed},
  stringstyle=\color{ForestGreen},
  showstringspaces=false,
  breaklines=true,
  breakatwhitespace=true,
  keywordstyle=\color{Purple},
  emph={def},
  emphstyle=\color{RoyalBlue}
}
\lstnewenvironment{itemlisting}[1][]
 {\mbox{}\vspace*{-\baselineskip}%
  \lstset{xleftmargin=\leftmargin, linewidth=\linewidth, #1}}
 {}

\makeatletter
\def\fps@figure{htbp}
\makeatother
\ifdefined\hasalgorithms
\fi

\newcommand{\docfooter}{\vfill\begin{tiny}Updated \today\end{tiny}}


\setupheader{\coursenumber\ -- \courseterm\ -- \coursedocname}
\pstitle{Problem Set 0: Setup and Mathematical Preliminaries}{\psZeroDue}


% Type an answer below each SOLUTION HERE marker. Increase the height argument
% on that studentanswer environment if additional space is required.
\newenvironment{studentanswer}[1]
  {\par\noindent\begin{minipage}[t][#1][t]{\linewidth}\ignorespaces}
  {\end{minipage}\par}

\begin{document}
\maketitle
\coursenotice
\begin{quote}\small\textbf{AI Policy.} You may use any AI tools you like to assist you with your assignments, and you should, as that is likely how you will work for the rest of your career. 
That being said, the written portions are practice for the in-class midterm and exam, which are closed book and on paper. You should strive to answer these yourself and use AI as a tutor and not an oracle.
\end{quote}

\callout{Please include an AI Tool Disclosure Below \vspace{50pt}}

\newpage
\section*{Overview}
This problem set has two goals. First, it reviews the mathematical foundations used throughout the course: linear algebra, basic probability, and derivatives. Second, it walks you through installing the course software environment and confirms that the simulator and testing tools run correctly on your machine.
This is intentionally the lightest problem set of the term. If any single step, and in particular the software install, takes you more than an hour, please stop and ask for help on \dgbf{Slack} rather than continuing to struggle alone.

\callout{\textbf{Review resources.} If a written problem below feels rusty,
we encourage you to review the underlying material before grinding through
the problem. The following free resources are matched to the parts of this
assignment: 3Blue1Brown's \emph{Essence of Linear Algebra} and \emph{Essence
of Calculus} video series on YouTube (Parts 1 and 3), Khan
Academy's probability unit (Part 2),
\href{https://immersivemath.com}{immersivemath.com} for interactive linear
algebra, and the free textbook \emph{Mathematics for Machine Learning}
(\href{https://mml-book.github.io}{mml-book.github.io}, Ch 2, 5, 6) as deeper reference.
PS0 is half problem set and half tutorial: the released solutions state each
rule before using it and link to further material.}

\vspace{10pt}
\textbf{What to submit.} Submit to \gradescopelink:\vspace{-10pt}
\begin{enumerate}\tightlist
\item \textbf{Written work (PDF)}: your answers to Part 1 and the C1
      trajectory image and observations. Space is
      provided below each problem; you may print this handout and write in
      it, annotate the PDF on a tablet, or use the provided
      \texttt{ps0/handout.tex}: enter each answer below its
      \texttt{\%\%\%\%\% SOLUTION HERE \%\%\%\%\%} marker, inside the
      corresponding \texttt{studentanswer} environment, and compile the file
      to produce your PDF. Assign each question to a page when you upload.
\item \textbf{Code}: your completed \texttt{warmup.py} from Part 2. Every
      test in this problem set is public, so if the autograder passes on
      your machine it will pass on submission.
\end{enumerate}

\newpage
\section*{Part 1: Written problems}

Work each problem on paper, without a calculator; every answer comes out in
small integers or simple fractions. Show your work in the space provided.

\subsection*{P1: Linear Algebra}

\begin{enumerate}[label=\alph*)]

\item Let $A = \begin{bmatrix} 2 & -1 \\ 0 & 3 \end{bmatrix}$ and $\mathbf{v} = \begin{bmatrix} 2 \\ 1 \end{bmatrix}$. Compute $A\mathbf{v}$.

\begin{studentanswer}{4cm}
%%%%% SOLUTION HERE %%%%%
\end{studentanswer}

\item The counterclockwise planar rotation matrix is
$R(\theta)=\begin{bmatrix}\cos\theta&-\sin\theta\\
\sin\theta&\cos\theta\end{bmatrix}$. Evaluate $B=R(90^\circ)$.
Let $A = \begin{bmatrix} 1 & 0 \\ 0 & -1 \end{bmatrix}$ be reflection
across the $x$-axis and let $\mathbf{v} = \begin{bmatrix} 1 \\ 2\end{bmatrix}$.
Compute $AB\mathbf{v}$ (rotate first, then reflect) and $BA\mathbf{v}$
(reflect first, then rotate). Do the two results agree? Based on this
example, is matrix multiplication commutative? Explain what this means about
the order in which transformations are applied.

\begin{studentanswer}{6cm}
%%%%% SOLUTION HERE %%%%%
\end{studentanswer}

\end{enumerate}

\newpage
\subsection*{P2: Probability}

\begin{enumerate}[label=\alph*)]
\item A discrete random variable $X$ takes value $0$ with probability $\tfrac{1}{2}$, value $2$ with probability $\tfrac{1}{4}$, and value $4$ with probability $\tfrac{1}{4}$. Compute $\E[X]$ and $\E[2X + 1]$.

\begin{studentanswer}{3.4cm}
%%%%% SOLUTION HERE %%%%%
\end{studentanswer}

\item A bump sensor beeps with probability $0.9$ when an obstacle is present and with probability $0.3$ when none is. The prior probability of an obstacle is $\tfrac{1}{3}$. The sensor beeps once. Use \href{https://en.wikipedia.org/wiki/Bayes_theorem}{Bayes' rule} to compute the posterior probability that an obstacle is present.

\begin{studentanswer}{3.4cm}
%%%%% SOLUTION HERE %%%%%
\end{studentanswer}

\item $X$ takes value $0$ with probability $\tfrac{1}{2}$ and value $2$
      with probability $\tfrac{1}{2}$. Compute $\E[X]$, $\E[X^2]$, and
      $\Var[X] = \E[X^2] - \E[X]^2$.

\begin{studentanswer}{3.2cm}
%%%%% SOLUTION HERE %%%%%
\end{studentanswer}

\item Two range sensors fail independently, each with probability
      $\tfrac{1}{4}$. Compute the probability that both fail, and the
      probability that at least one fails. Which of the two computations
      \emph{required} the independence assumption?

\begin{studentanswer}{3.4cm}
%%%%% SOLUTION HERE %%%%%
\end{studentanswer}

\end{enumerate}

\newpage
\subsection*{P3: Calculus}

\begin{enumerate}[label=\alph*)]
\item For $f(x, y) = x^2 + 3xy$, compute the gradient
$\nabla f = \left(\tfrac{\partial f}{\partial x},
\tfrac{\partial f}{\partial y}\right)$ and evaluate it at $(1, 2)$.
Then check each component numerically with the central-difference formula
\[
\frac{\partial f}{\partial x_i}(\mathbf{x}) \approx
\frac{f(\mathbf{x}+\epsilon\mathbf{e}_i)-
f(\mathbf{x}-\epsilon\mathbf{e}_i)}{2\epsilon}
\]
using $\epsilon=0.1$.
Compute both numerical estimates and state whether they agree with your
analytical gradient.

\begin{studentanswer}{4cm}
%%%%% SOLUTION HERE %%%%%
\end{studentanswer}

\item Use the chain rule: for $g(x) = (2x + 1)^2$, compute $g'(x)$ and evaluate $g'(1)$.

\begin{studentanswer}{4cm}
%%%%% SOLUTION HERE %%%%%
\end{studentanswer}

\item Consider the loss function $L(w) = \tfrac{1}{2}(wx - y)^2$ with the single data point $x = 3$, $y = 5$. This is a standard squared-error loss in machine learning. Compute $\tfrac{dL}{dw}$, evaluate it at $w = 2$, and take one gradient descent step $w_1 = w_0 - \lr \tfrac{dL}{dw}(w_0)$ with step size $\lr = 0.1$.

\begin{studentanswer}{5cm}
%%%%% SOLUTION HERE %%%%%
\end{studentanswer}

\end{enumerate}

\newpage
\subsection*{P4: Deriving the Normal Equations \gradstar}

The least squares algorithm can be derived directly from the two toolboxes this problem set reviews. Let $A$ be an $m \times n$ matrix of inputs and $\mathbf{b}$ an $m$-vector of targets, and define $f(\mathbf{x}) = \|A\mathbf{x} - \mathbf{b}\|^2$.

\begin{enumerate}[label=\alph*)]
\item Show that any minimizer satisfies the normal equations $A^\top A\,\mathbf{x} = A^\top \mathbf{b}$.

\begin{studentanswer}{8cm}
%%%%% SOLUTION HERE %%%%%
\end{studentanswer}

\item In one or two sentences, when is the solution unique?

\begin{studentanswer}{3cm}
%%%%% SOLUTION HERE %%%%%
\end{studentanswer}

\end{enumerate}

\newpage
\section*{Part 2: Programming}

The programming part has three steps: install the course software
environment, run a provided script to confirm the environment works, and
implement five short functions that mirror the written problems above.

\subsection*{Installing the course environment}

The course code requires \textbf{Python 3.12} and a terminal. The
instructions below assume no prior experience with either; if you have
set up Python projects before, the summary is: clone the repository, create a
virtual environment at its root, and \texttt{pip install -r requirements.txt}.

\textbf{Step 1: get the course code.} The graphical option is GitHub Desktop:
choose \textbf{File $>$ Clone repository}, select
\url{https://github.com/plancherb1/cosc-281-f26}, and choose where to keep the
course folder. From a terminal, the equivalent commands are:

\begin{lstlisting}[language={}]
git clone git@github.com:plancherb1/cosc-281-f26.git
cd cosc-281-f26
\end{lstlisting}

Clone only once. When a new problem set is announced, use
\textbf{Repository $>$ Pull} in GitHub Desktop or run \texttt{git pull} from
the course folder. A public fork is not required and should not contain your
answers.

\textbf{Step 2: open a terminal and check Python and Tk.} Tk is the GUI
component that opens the driving window; \texttt{pip} cannot install it.

\begin{itemize}\tightlist
\item \textbf{macOS}: press Cmd-Space, type \texttt{Terminal}, and press
      return. Check your Python version with \texttt{python3 --version}.
      If the version is below 3.12 or the command is not found, install
      Python 3.12 from \url{https://www.python.org/downloads/} (use the
      macOS universal installer, then close and reopen the terminal). Run
      \texttt{python3 -m tkinter}; a small window should open. Close it after
      the check.
\Needspace{8\baselineskip}
\item \textbf{Windows}: install Python 3.12 from
      \url{https://www.python.org/downloads/}. On the first installer
      screen, \textbf{check the box labeled ``Add python.exe to PATH''}
      before clicking Install; this is the step most people miss. Then
      open PowerShell (press the Windows key, type \texttt{PowerShell},
      press enter) and check \texttt{python --version}. Note that on
      Windows the command is \texttt{python}, not \texttt{python3}. Then run
      \texttt{python -m tkinter}. If no small window opens, modify the Python
      installation and enable \textbf{Tcl/Tk and IDLE}.
\item \textbf{Linux}: open your terminal application and check
      \texttt{python3 --version}. On Ubuntu or Debian, install what you
      need with \texttt{sudo apt install python3 python3-venv
      python3-tk}. Run \texttt{python3 -m tkinter} and close the small window
      after it appears.
\end{itemize}

\textbf{Step 3: navigate to the course folder.} If you used GitHub Desktop,
open its \textbf{Repository} menu and choose \textbf{Open in Terminal}. You
can also type \texttt{cd} followed by a space and the path to the cloned
folder. On most systems you can type \texttt{cd }
(with the trailing space) and then drag the folder from your file browser
into the terminal window to fill in the path. Confirm you are in the right
place: \texttt{ls} (macOS/Linux) or \texttt{dir} (Windows) should list
\texttt{sim}, \texttt{ps0}, and \texttt{requirements.txt}.

\Needspace{12\baselineskip}
\textbf{Step 4: create and activate a virtual environment.} A virtual
environment is a private copy of Python's package installation for this
course, so nothing here conflicts with anything else on your machine.

\begin{itemize}\tightlist
\item \textbf{macOS / Linux}:
\begin{lstlisting}[language={}]
python3 -m venv --prompt COSC281 .venv
source .venv/bin/activate
\end{lstlisting}
\item \textbf{Windows (PowerShell)}:
\begin{lstlisting}[language={}]
python -m venv --prompt COSC281 .venv
.venv\Scripts\Activate.ps1
\end{lstlisting}
      If PowerShell refuses to run the activation script with an error
      about execution policies, run
      \texttt{Set-ExecutionPolicy -ExecutionPolicy RemoteSigned -Scope
      CurrentUser} once, answer yes, and try the activation again.
\end{itemize}

When activation succeeds, your terminal prompt begins with
\texttt{(COSC281)}. This confirms that the course environment, rather than
another project's environment, is active. You must activate it (the second
command only) in every new terminal window you use for course work.

\textbf{Step 5: install the course packages.}

\begin{lstlisting}[language={}]
pip install -r requirements.txt
\end{lstlisting}

This downloads and installs the pinned course packages, including the
course simulator (\texttt{coursesim}, which lives in \texttt{sim/}; the
course scripts find it on their own, and there is nothing separate to
install). It may take a few minutes.

\textbf{Step 6: verify.} Run the public test suite:

\begin{lstlisting}[language={}]
python ps0/code/autograder.py
\end{lstlisting}

The three environment tests should pass immediately. The ten
\texttt{warmup} tests will fail with \texttt{NotImplementedError} until
you complete C2 below; that is expected at this stage.

If anything above fails, copy the \emph{entire} error message and post it
on Slack along with your operating system. Installation problems are
normal, they are fastest to fix with the full error text, and helping you
through them is exactly what week one is for.

\Needspace{15\baselineskip}
\subsection*{C1: Drive the car and plot the trajectory}

\textbf{Where to answer: your written-work PDF.} If you use the editable
TeX handout, place the image and observations inside the C1
\texttt{studentanswer} environment.

With the environment activated, run the teleoperation script:

\begin{lstlisting}[language={}]
python ps0/code/teleop_drive.py --map sim/maps/open.yaml \
    --out ps0/figures/ps0_teleop.png
\end{lstlisting}

A window opens showing a small car on an open map. Drive it with the arrow
keys (or WASD) for at least 20 seconds of simulated time, then press
\texttt{q} to quit. The script saves \texttt{ps0\_teleop.png}, an image of
the trajectory you drove drawn over the map. This script requires a
display; if you work over ssh, use \texttt{ssh -X} or run locally. If no
window opens, do not use a scripted trajectory as your answer; follow the
installation fix printed by the command and retry until you can drive
interactively.

Include the image in your PDF along with two or three sentences of
observation: what motions could the car execute and what motions could it
not (try to turn in place), and one thing about its behavior that you did
not expect. The purpose of this task is to confirm that your environment
works end to end and to give you direct experience with the car's steering
constraint.

If you use the editable TeX handout, the image can be inserted at the C1
marker with:
\begin{lstlisting}[language={}]
\includegraphics[width=0.8\linewidth]{ps0/figures/ps0_teleop.png}
\end{lstlisting}
Write your two or three sentences below the image command.

\begin{studentanswer}{8cm}
%%%%% SOLUTION HERE %%%%%
\end{studentanswer}

\Needspace{24\baselineskip}
\subsection*{C2: Five warmup functions}

\textbf{Where to answer: \texttt{warmup.py}.} Write these implementations in
the provided Python file and upload that file as the code submission. Do not
paste the implementations into the written-work PDF.

The file \texttt{ps0/code/warmup.py} contains five function
stubs. Each is the computational counterpart of a specific written problem.
The docstring above each stub gives its contract, formula, and matching
paper example. Read that contract before writing the function.

\begin{itemize}\tightlist
\item \texttt{matvec(A, v)} implements P1(a). Use two explicit nested loops.
      You may use NumPy to convert the inputs, inspect their shapes, and
      allocate the output, but multiplication and accumulation must be scalar
      operations inside your loops. Do not use \texttt{@}, \texttt{np.dot},
      \texttt{np.matmul}, \texttt{np.einsum}, or another routine that computes
      the product for you.
\item \texttt{rot2d(theta)} constructs the rotation matrix from P1(b), with
      \texttt{theta} expressed in radians.
\item \texttt{expectation(values, probs)} generalizes the sum in P2(a).
\item \texttt{bayes\_posterior(...)} generalizes the update in P2(b).
\item \texttt{central\_diff\_grad(f, x, eps)} implements the numerical
      check in P3(a), one coordinate at a time.
\end{itemize}

Check your work by re-running the autograder:

\begin{lstlisting}[language={}]
python ps0/code/autograder.py
\end{lstlisting}

All thirteen tests pass when the environment and your five functions are
correct. Every test in this problem set is public: there are no hidden
tests, so a passing run on your machine is a passing run on Gradescope.
Submit \texttt{warmup.py} (the only file you edit) along with your PDF.

\docfooter
\end{document}
